\section{Objetivos}
\subsection{Objetivo general}
Sintetizar y caracterizar puntos cuánticos de carbono dopados con Nitrógeno(N-CQDs) mediante pirólisis asistida por microondas, evaluando el efecto del tipo de precursor, la estequiometría, el tiempo de reacción y el medio químico sobre sus propiedades ópticas y microestructurales, e implementar una plataforma metrológica de fotometría digital RGB para la cuantificación directa de su emisión fotoluminiscente.

\subsection{Objetivos específicos}
\begin{enumerate}
    \item Evaluar el efecto de la relación molar (1:6 y 1:10) y del tiempo de pirólisis en la ruta sintética base de ácido cítrico y urea, diferenciando la formación de fases masivas no luminiscentes frente a suspensiones coloidales fluorescentes.
    
    \item Desarrollar una ruta sintética verde utilizando clara de huevo (ovoalbúmina) como bioprecursor de nitrógeno y azufre, determinando la ventana cinética óptima (\textit{sweet spot}) a 1:15~minutos de pirólisis y calculando su \textit{bandgap} óptico directo ($E_g = 3{,}33$~eV).
    
    \item Analizar la influencia del medio ácido mediante la adición de vinagre en la reacción con ácido cítrico y urea, evaluando el corrimiento batocrómico resultante en los espectros de absorción UV-Vis y su impacto en la composición elemental por EDX.
    
    \item Implementar un sistema de fotometría digital RGB en cámara oscura con procesamiento de imágenes en Python, estableciendo la luminancia neta ($Y_{\text{neta}}$) para la cinética de reacción y la relación ratiométrica verde-azul ($G/B$) como indicador del corrimiento espectral.
\end{enumerate}